% ============================================================================
% THE WITNESS (paper #2) -- v0.1 skeleton, opened 2026-08-02.
%
% Authorized by the C2 ruling (docs/exploration/council-0802.md) and the two
% christenings in docs/blueprints/flags.md: THE PEARL
% (logChowla2_witnessed_scale_flat_L_v2 and its lineage) and THE TOLL (the
% entropy-decrement width law of the flat tower).
%
% TODO(title): the working title below is a placeholder chosen only so that
% both named theorems appear in it. The title is the author's to settle.
%
% House conventions, inherited verbatim from papers/flagship/main.tex:
%   * amsart, same preamble;
%   * every stated result carries its Lean declaration name via \leanname{};
%   * Appendix A is the statement -> declaration -> axiom audit table;
%   * voice: problem-first; every strong adjective paid by a number; every
%     conditional named at the statement, never in a footnote;
%   * author/thanks block and Acknowledgments as in the flagship.
% Two words are banned from the prose of this paper by house rule: "honest"
% (and "honestly") and "load-bearing". Em-dashes are permitted.
%
% SECTION STATUS at v0.1
%   1  Introduction ........ DRAFTED (conservative; every hypothesis named)
%   2  The Toll ............ DRAFTED IN FULL (the chapter this draft exists for)
%   3  The flat tower ...... DRAFTED IN FULL (the technical heart of section 2)
%   4  The Pearl ........... STUB -- awaits BAND-WINDOW's final hypothesis list
%   5  The classical limit . HALF -- containment + hoist drafted; final form TODO
%   6  Method .............. STUB -- contents listed, prose not written
%   A  Statement index ..... Toll rows FILLED; Pearl rows stubbed
%
% No local LaTeX toolchain on the build machine: this file has never been
% compiled. Compile checks happen author-side.
% ============================================================================
\documentclass[11pt]{amsart}
\usepackage[margin=1.1in]{geometry}
\usepackage{amsmath,amssymb,amsthm}
\usepackage{booktabs}
\usepackage[colorlinks=true,linkcolor=blue!60!black,citecolor=blue!60!black,urlcolor=blue!60!black]{hyperref}
\usepackage{microtype}
\usepackage{xcolor}

\newcommand{\leanname}[1]{\par\noindent{\small\texttt{#1}}}
\newcommand{\leaninline}[1]{{\small\texttt{#1}}}

\newtheorem{theorem}{Theorem}[section]
\newtheorem{corollary}[theorem]{Corollary}
\newtheorem{definition}[theorem]{Definition}
\newtheorem{proposition}[theorem]{Proposition}
\newtheorem{lemma}[theorem]{Lemma}
\theoremstyle{remark}
\newtheorem{remark}[theorem]{Remark}

\title[The Pearl and the Toll]{The Pearl and the Toll: a Witnessed-Scale
Chowla Regime and the Width Law of the Entropy Decrement}
\author{Jason Hickey}
\thanks{Developed in collaboration with Claude (Anthropic's Fable~5 and
Opus models); see \S\ref{sec:method} and the repository's commit ledger,
in which every commit is co-authored. Written in a personal capacity.}
\date{August 2, 2026}

\begin{document}

\begin{abstract}
Every theorem in this paper is verified by the Lean~4 kernel over
\textsf{mathlib}, with axiom base exactly
$\{\texttt{propext}, \texttt{Classical.choice}, \texttt{Quot.sound}\}$;
nothing rests on informal argument. The mathematics---the designs, the
proofs, and their formalization---is joint work of the author and Claude
(Anthropic's Fable~5 and Opus models), carried out under the author's
direction and ratification; the commit ledger records the collaboration
line by line.

We prove two theorems about the entropy-decrement route to the
logarithmically averaged two-point Chowla statement. The first, the
\emph{toll}, is unconditional: along the flat tower of scales
$H_{j+1} = H_j\lfloor 2A\log H_j\rfloor$ at design constant $A \ge 1$, the
telescoped entropy decrement equals a potential difference to within
$5\%$, and consequently a tower long enough to exhaust the one-bit entropy
budget $\log 2$ must multiply the doubly-logarithmic width
$\log 2A + \log\log H$ by a factor in the bracket
$[\,4^{(20/21)A},\ (21/20)\,4^{A}\,]$ --- exponent pinned within $4.8\%$,
constant within $5\%$, with the continuum value $4^{A}$ inside. Crossing
the entropy budget is thus a conservation law rather than an optimization:
the price is fixed, and a design chooses only the coordinate in which to
pay it. The second theorem, the \emph{pearl}, is conditional and we name
its conditionality at the statement: at every design constant $A$ above an
explicit floor there is a Chowla regime whose window base is pinned at
$\lceil e^{e^{3.2A}}\rceil$, and for that regime the logarithmically
averaged two-point Chowla correlation does not fail at the witnessed
scale, on a list of hypotheses given in full in \S\ref{sec:pearl}: one
analytic constant of the band lane, eight numeric riders on constants the
theorem itself produces, and two carried predicates. A corollary removes
the two Siegel-genre riders of that list at the price of effectivity, by
the classical limit; making that limit available required a quantifier hoist
that the formal statement, and only the formal statement, made visible.
\end{abstract}

\maketitle

% ============================================================================
\section{Introduction}\label{sec:intro}
% ============================================================================

The logarithmically averaged two-point Chowla statement,
\begin{equation}\label{eq:chowla}
\frac{1}{\log x} \sum_{n \le x} \frac{\lambda(n)\lambda(n+1)}{n}
\;\longrightarrow\; 0 ,
\end{equation}
is a theorem of Tao \cite{TaoChowla}, proved by an entropy decrement: if
the correlation fails at a scale, the per-symbol entropy of the Liouville
window drops by a definite amount when the window length is multiplied,
and the drops telescope. Entropy per symbol of a $\pm1$-valued window lies
in $[0, \log 2]$, so the total available decrement over any tower of
scales is one bit. The proof therefore has a budget, and the budget is
spent by a tower of scales: build the tower long enough that the summed
decrement exceeds $\log 2$, and the failure hypothesis dies.

What does that tower cost? The question is not rhetorical in a formal
development, where the tower is an actual sequence of natural numbers, its
length an actual index, and its terminal scale an actual bound that
downstream lemmas must accommodate. In the campaign behind this paper the
answer was, three separate times, the difference between a route that
closed and a route that did not. The first of our two theorems answers it
exactly.

\subsection{The toll}
Fix a real design constant $A \ge 1$ and an integer base $B$, and let the
\emph{flat tower} be $H_0 = B$, $H_{j+1} = H_j\lfloor 2A\log H_j\rfloor$;
its telescoped decrement at length $J$ is
$S_J = \sum_{j<J} \bigl(2A\log H_j\bigr)^{-1}$. Write $L = \log H$,
$\lambda = \log L$, $c = \log 2A$.

\begin{theorem}[The toll; \S\ref{sec:toll}, unconditional]\label{thm:toll-intro}
Let $A \ge 1$, $B \ge 4\cdot 10^{6}$, and let the base satisfy the flat
floor $100(c + \lambda_0) \le L_0$. Then any tower length $J$ at which the
decrement crosses the one-bit budget, $\log 2 < S_J$, forces
\[
c + \lambda_J \;\ge\; 4^{(20/21)A}\,\bigl(c + \lambda_0\bigr),
\]
and at the least such length the width does not overshoot:
\[
c + \lambda_{J} \;\le\; \tfrac{21}{20}\,4^{A}\,\bigl(c + \lambda_0\bigr).
\]
\leanname{Salt.Entropy.Chowla.towerFlat\_width\_ge,
Salt.Entropy.Chowla.towerFlat\_width\_le}
\end{theorem}

The two halves bracket the width multiplier in
$[\,4^{(20/21)A},\ (21/20)\,4^{A}\,]$, which contains the continuum value
$4^{A}$: the exponent is pinned to within $20/21 = 0.952\ldots$ of the
truth and the constant to within $21/20$. Both halves come from one law
(Theorem~\ref{thm:master}): the decrement $S_J$ \emph{is} a potential
difference, $S_J = (2A)^{-1}(w_J - w_0)$ with $w = \log(c+\lambda)$, to
within a factor $21/20$. The whole content of the floor function in the
tower's definition is the $20/21$.

We read Theorem~\ref{thm:toll-intro} as a conservation law. Crossing the
entropy budget costs a fixed amount, $4^{A}$ up to $5\%$ in each direction,
and no design decision changes the amount; a design decides only which
coordinate is billed. The flat threshold $\kappa = H/(A\log H)$ pays in
doubly-logarithmic width, with $A$ a constant of the designer's choosing.
The threshold of the original argument, $\kappa = H/(\log H\log\log\log H)$,
pays the same invoice with $\log\log\log H$ in the role of $A$ --- and
since that is not a constant, the bill arrives on the scale axis instead,
as the triple-exponential floor $\lceil e^{e^{e^{X}}}\rceil$ that the flat
road replaces by the single exponential
$\lceil e^{\max(4X,\,2\log A + 2)}\rceil$
(\leaninline{budgetFloor} against \leaninline{budgetFloorFlat}). \S\ref{sec:toll}
makes this two-axis reading precise and states exactly which part of it is
kernel-checked. The width half is not special to a constant threshold: for an
arbitrary threshold schedule $\varphi_j \ge A_b$ driving the same step, crossing
the budget still forces a factor at least $4^{(20/21)A_b}$, and the constant
schedule at $A_b$ pays at most $\frac{21}{20}4^{A_b}$
(Theorem~\ref{thm:shapetoll} and Corollary~\ref{cor:extremal}).

\subsection{The pearl}
The second theorem is conditional, and we state its conditionality here
rather than later. In the form the corpus carries it: for every $A_0 \ge
162$ there exist constants
$(\varepsilon, C_g, K_c, \delta_0, C_t, A, \beta, x_0, H_{\mathrm{opq}},
M_{\mathrm{fl}}, C_q, c_s, T_0, K_q, K_s, C)$ with $A \ge \max(162, A_0)$
such that, \emph{given} eight numeric hypotheses on those constants and
two carried predicates, there is a Chowla regime $R$ with window base
pinned at the design base, terminal scale beyond any prescribed function
of the regime's own data, and
$\lnot\,\mathrm{logChowla2Fails}(R.\varepsilon, R.x, R.\omega)$ --- that
is, the logarithmically averaged two-point correlation at the witnessed
window does not exceed $\varepsilon\log\omega$ in absolute value.
\leanname{Salt.MR.logChowla2\_witnessed\_scale\_flat\_L\_v2}

The eight hypotheses and two predicates are listed in full in
\S\ref{sec:pearl}; none is hidden, and the paper claims nothing
unconditional about \eqref{eq:chowla}. Two of the eight are of Siegel
genre --- windows on constants that the theory does not effectively bound
--- and they are removed by the corollary of \S\ref{sec:limit}:
\leanname{Salt.MR.logChowla2\_ineffective}
which carries the band-lane hypothesis, six numeric riders
($K_c \le 2^{539}$, $C_t \le 2^{23}$, $e^{-100} \le c_s$,
$T_0 \le e^{e^{100}}$, $K_q \le e^{100}$, $e^{-100} \le K_s$) and the two
carried predicates, and no Siegel item of any kind.

\subsection{Effective, ineffective, and the order of quantifiers}
The two statements stand in the anatomy familiar from the theory of the
exceptional zero: an \emph{effective} statement conditional on a bound
relating a fixed constant to a growing window contains an
\emph{unconditional-in-that-constant} statement whose own constant is then
ineffective. Here the growing windows are double exponentials in the
design constant $A$ ($x_0 \le e^{e^{3.2A}/10}$ and
$H_{\mathrm{opq}} \le \lceil e^{e^{3.2A}}\rceil$), the constants they
window are fixed, and choosing $A$ above both discharges them; what is
lost is effectivity, since the choice is made by
\texttt{Classical.choice}. This is exactly the caveat the analytic
literature attaches to Siegel's theorem.

The formalization contributed one thing to that argument which prose would
not have surfaced: at the quantifier order in which the conditional
theorem was first proved, the classical limit is \emph{not available}. The
statement reads $\forall A_0 \ge 162,\ \exists (\ldots, x_0,
H_{\mathrm{opq}}, A, \ldots)$, so every constant is chosen after the
caller's $A_0$ and nothing in the statement forbids $x_0$ and
$H_{\mathrm{opq}}$ from moving when $A_0$ does; two applications of
\texttt{Classical.choose} at two different $A_0$ are choices from two
different propositions, and proof irrelevance does not bridge them. The
constants are in fact $A_0$-free at the level of the proof, and making
that a fact of the \emph{statement} required re-cutting seven theorems
along the road with the quantifier hoisted (\S\ref{sec:limit}). We record
this as a methodological finding: a formal statement carries its
quantifier order as data, and an argument the prose would perform without
comment can be blocked by it.

\subsection{What is in this paper}
\S\ref{sec:toll} states and proves the toll and its exports.
\S\ref{sec:flat} develops the flat tower --- the construction, the floor
invariant, the per-step bracket --- which is the technical content of
\S\ref{sec:toll}'s proof. \S\ref{sec:pearl} states the pearl with its
complete hypothesis list. \S\ref{sec:limit} gives the classical limit and
the quantifier hoist. \S\ref{sec:method} describes the method, which is
treated at length in a companion paper \cite{companion} and in the
program's first paper \cite{fulcrum}. Appendix~\ref{app:index} maps every
numbered statement to its Lean declaration and its axiom audit
(Table~\ref{tab:index}).

Conditionality is tabulated in one place, Table~\ref{tab:conditional};
every conditional statement in the paper names its exact residual
hypothesis at the point of statement.

% ============================================================================
\section{The toll}\label{sec:toll}
% ============================================================================

\subsection{The budget and the tower}
The entropy-decrement argument runs on a tower of window lengths. At
window length $H$ let
$e(H) := \mathbb{H}\bigl[\lambda\text{-window of length } H\bigr]/H$
be the per-symbol entropy of the Liouville window under the logarithmic
measure of the regime. Two endpoint facts fix the budget: $e(H) \le \log
2$ for every $H \ge 1$, since the window is $\pm1$-valued
(\leaninline{Salt.Entropy.Chowla.entropy\_per\_symbol\_le}), and $e(H) \ge
0$ (\leaninline{Salt.Entropy.Chowla.entropy\_nonneg\_per\_symbol}). The
decrement step says that if the correlation fails at level $j$ then
$e(H_{j+1}) - e(H_j) \le -\,1/(2A\log H_j)$ at the flat threshold
$\kappa = H/(A\log H)$; summing over $j < J$ gives
\[
e(H_J) \;\le\; e(H_0) - S_J, \qquad
S_J \;:=\; \sum_{j<J} \frac{1}{2A\log H_j},
\]
\leaninline{Salt.Entropy.Chowla.tower\_telescopeFlat}, and with the two
endpoint facts, $S_J \le \log 2$. So the argument closes exactly when the
tower is long enough that $S_J > \log 2$: one bit, in nats, is the entire
capital of the method.

The tower cannot be extended for free, because $H_j$ grows and the
summands $1/(2A\log H_j)$ shrink. The question the toll answers is what
the terminal scale must be.

Convention for \S\ref{sec:toll} and \S\ref{sec:flat}: every declaration
cited below lives in \texttt{Salt/Entropy/Chowla/TowerFlat.lean} of the
repository unless a different file is named at the citation, and all are
in the namespace \leaninline{Salt.Entropy.Chowla}.

\begin{definition}[The flat tower and its decrement]\label{def:flattower}
For a real $A$ and naturals $a$, $B$, set $H_0 = a\,B$ and
\[
H_{j+1} \;=\; H_j \cdot \bigl\lfloor 2A \log H_j \bigr\rfloor ,
\]
and $S_J := \sum_{j<J} \bigl(2A\log H_j\bigr)^{-1}$. Throughout this paper
$a = 1$, so $H_0 = B$.
\leanname{Salt.Entropy.Chowla.chowlaTowerFlat,
Salt.Entropy.Chowla.towerDropSumFlat}
\end{definition}

The multiplier $\lfloor 2A\log H_j\rfloor$ is chosen so that the tower
funds its own threshold: the step's two error terms fit inside half the
per-step drop $1/(2A\log H_j)$ for every $\varepsilon \le 1/2$, $A \ge 1$
and $\log H \ge 15$
(\leaninline{Salt.Entropy.Chowla.flat\_hkey}). The name \emph{flat} marks
the contrast with the original threshold $H/(\log H \log\log\log H)$: the
third logarithm has been replaced by a constant the designer picks.

\begin{definition}[The potential]\label{def:potential}
Write $L_j = \log H_j$, $\lambda_j = \log L_j$, $c = \log 2A$, and
\[
w_j \;=\; \log\bigl(c + \lambda_j\bigr).
\]
\leanname{Salt.Entropy.Chowla.towerLF, towerLamF, flatC, towerWF}
\end{definition}

The potential is forced by the arithmetic of the step, not chosen for
convenience. One step adds $d_j = \log\lfloor 2A L_j\rfloor$ to $L$, and
$d_j = c + \lambda_j$ exactly up to the floor function; so
$\Delta\lambda_j \approx (c+\lambda_j)/L_j$ and
$\Delta w_j \approx 1/L_j$. That is: the summand $1/(2AL_j)$ of $S_J$
\emph{is} $\Delta w_j/(2A)$. The next theorem says so with explicit
two-sided constants.

\subsection{Hypotheses}
Three hypotheses are in force for every statement of this section, and
each is named at every statement that uses it.
\begin{itemize}
\item[(H1)] $A \ge 1$: the design constant is at least one.
\item[(H2)] $B \ge 4\cdot 10^{6}$: the house regime floor, which supplies
  $\log B \ge 15$ and, along the tower, $L_j \ge 15$ and $\lambda_j \ge 2$
  at every level.
\item[(H3)] the \emph{flat floor} $100\,(c + \lambda_0) \le L_0$, i.e.
  $100\,(\log 2A + \log\log B) \le \log B$.
\end{itemize}
(H3) is what makes the second-order factors of the step $1/100$-small, and
it propagates up the tower: $x \mapsto x - 100(c + \log x)$ is increasing
past $x = 100$, so the invariant survives every step
(\leaninline{Salt.Entropy.Chowla.towerFlat\_inv}). It is a mild demand. At
$A = 10^{30}$, where $c = 69.7\ldots$, it asks $\log B \ge 8\cdot10^{3}$
(a numerical remark; the formal statement is the displayed inequality). In
the use made of the tower in \S\ref{sec:pearl}, where the design law
$3.2A \le \log\log B$ holds with $A \ge 26$, the flat floor is not a
hypothesis at all but a theorem:
\leanname{Salt.Entropy.Chowla.flatFloor\_of\_design}

\subsection{The master law}

\begin{theorem}[The master law: the decrement is a potential
difference]\label{thm:master}
Assume (H1), (H2), (H3). Then for every $J$,
\[
\frac{1}{2A}\bigl(w_J - w_0\bigr)
\;\le\; S_J \;\le\;
\frac{21}{20}\cdot\frac{1}{2A}\bigl(w_J - w_0\bigr).
\]
\leanname{Salt.Entropy.Chowla.towerDropSumFlat\_ge\_log\_ratio,
Salt.Entropy.Chowla.towerDropSumFlat\_le\_log\_ratio\_mul}
\end{theorem}

The two-sided constant is $21/20$: the telescoped decrement of the flat
tower is the potential difference $w_J - w_0$, divided by $2A$, to within
$5\%$. The lower half is exact --- no relative loss at all --- because both
prices charged by the floor function point the same way in that direction;
the $1/20$ of the upper half pays three items, two second-order factors of
the step ($L' \le (101/100)L$ and $c + \lambda' \le (c+\lambda) + 1/100$,
both from (H3)) and the floor price $1/500$.

\subsection{The toll}

\begin{theorem}[The toll, necessity]\label{thm:tollge}
Assume (H1), (H2), (H3). For \emph{every} tower length $J$ at which
$\log 2 < S_J$,
\[
\bigl(c + \lambda_0\bigr)\, 4^{(20/21)A} \;\le\; c + \lambda_J .
\]
\leanname{Salt.Entropy.Chowla.towerFlat\_width\_ge}
\end{theorem}

\begin{theorem}[The toll, no overshoot]\label{thm:tollle}
Assume (H1), (H2), (H3), and let $J_{\min}$ be the least $J$ with
$\log 2 < S_J$ (it exists, by Proposition~\ref{prop:crossing}). Then
\[
c + \lambda_{J_{\min}} \;\le\; \frac{21}{20}\,4^{A}\,\bigl(c + \lambda_0\bigr).
\]
\leanname{Salt.Entropy.Chowla.towerFlat\_width\_le}
\end{theorem}

\begin{proposition}[The crossing is reached, at a named
level]\label{prop:crossing}
Assume (H1), (H2), (H3), and set
$J^{\star} := \bigl\lceil \exp\bigl(2(c+\lambda_0)4^{A}\bigr)\bigr\rceil$.
Then $\log 2 < S_{J^{\star}}$; in particular the set of crossing lengths
is nonempty and $J_{\min}$ is well defined and positive.
\leanname{Salt.Entropy.Chowla.towerFlatJstar,
towerFlat\_crossing\_at\_Jstar, towerFlat\_crossing\_exists,
towerFlatJmin, towerFlatJmin\_spec, towerFlatJmin\_pos}
\end{proposition}

Together, Theorems~\ref{thm:tollge} and \ref{thm:tollle} bracket the width
multiplier of a minimal crossing tower:
\begin{equation}\label{eq:bracket}
4^{(20/21)A} \;\le\; \frac{c+\lambda_{J_{\min}}}{c+\lambda_0}
\;\le\; \frac{21}{20}\,4^{A},
\end{equation}
and the continuum value $4^{A}$ --- what the same computation gives with
the floor function and the discreteness removed --- lies inside. The
bracket's own width is a factor $4^{A/21}\cdot(21/20)$: the exponent base
is determined, the exponent to within $4.8\%$, the constant to within
$5\%$.

To see the size of the toll, take the working point of \S\ref{sec:pearl},
$A = 162$, where the base is pinned so that $\lambda_0 \ge 3.2A = 518.4$
and $c = \log 324 = 5.78\ldots$. Then \eqref{eq:bracket} puts
$\lambda_{J_{\min}}$ between $4\cdot10^{95}$ and $1.9\cdot10^{100}$: the
terminal scale of the tower is $H_{J_{\min}} = \exp\exp(\lambda_{J_{\min}})$,
a doubly exponential object whose doubly-logarithmic width has been
multiplied by about $10^{93}$. (Numerical illustration; the formal content
is the displayed inequalities.)

\begin{corollary}[The $A$-free export]\label{cor:export}
Let $A \ge 26$, $B \ge 4\cdot10^{6}$, and let the design law
$3.2A \le \log\log B$ hold. Then, at the minimal crossing length,
\[
\log\log H_{J_{\min}} \;\le\; \exp\!\left(\frac{\log\log B}{2}\right).
\]
\leanname{Salt.Entropy.Chowla.towerFlat\_width\_export}
\end{corollary}

Corollary~\ref{cor:export} is Theorem~\ref{thm:tollle} with the design law
substituted and $A$ eliminated: under $3.2A \le \lambda_0$ the whole right
side $\frac{21}{20}4^{A}(c+\lambda_0)$ is below $e^{\lambda_0/2}$, with the
step from the design point $\lambda_0 = 3.2A$ to the whole half-line
$\lambda_0 \ge 3.2A$ carried by a single application of $1 + x \le e^{x}$
(\leaninline{Salt.Entropy.Chowla.flatWidth\_coef},
\leaninline{flatWidth\_engine\_at}, \leaninline{flatWidth\_engine}). The
$A$-free form is what \S\ref{sec:pearl}'s road consumes; stating it
$A$-free is what keeps the design constant out of a hundred downstream
statements.

\subsection{The proof}
All four statements come from one per-step bracket, proved in
\S\ref{sec:flat} as Lemma~\ref{lem:step}: under the invariant (H3) at
level $j$,
\begin{equation}\label{eq:stepbracket}
\frac{20/21}{L_j} \;\le\; w_{j+1} - w_j \;\le\; \frac{1}{L_j}.
\end{equation}
Given \eqref{eq:stepbracket} the rest is bookkeeping in three moves.

\emph{Telescoping.} Multiply \eqref{eq:stepbracket} by $1/(2A)$ and
compare with the summand $1/(2AL_j)$ of $S_J$ term by term; summing over
$j < J$ and telescoping $\sum_{j<J}(w_{j+1}-w_j) = w_J - w_0$ gives
Theorem~\ref{thm:master}, the upper half acquiring the factor $21/20 =
(20/21)^{-1}$ and the lower half none.

\emph{Exponentiating, necessity.} If $\log 2 < S_J$, the upper half of the
master law gives $\log 2 < \frac{21}{40A}(w_J - w_0)$, i.e.
$w_J - w_0 > \frac{40}{21}A\log 2 = \frac{20}{21}A\log 4$. Since
$w = \log(c+\lambda)$, exponentiating is Theorem~\ref{thm:tollge}. The
$\log 4$ is where the base $4$ of the toll comes from: the budget is
$\log 2$ and the drop carries a factor $\frac{1}{2A}$, so the exponent
carries $2\log 2$.

\emph{Exponentiating, no overshoot.} At $J_{\min} = k+1$, minimality gives
$S_k \le \log 2$, and the lower half of the master law converts that into
$w_k - w_0 \le 2A\log 2 = A \log 4$. One further step costs at most
$1/L_k \le 1/300 \le 1/21 \le \log(21/20)$, using $L_k \ge 300$ from (H3)
and $c + \lambda_k \ge 3$. Hence
$w_{J_{\min}} - w_0 \le \log\frac{21}{20} + A\log 4$, which is
Theorem~\ref{thm:tollle}. Proposition~\ref{prop:crossing} is the crude
direction: each step adds at least $1$ to $L$, so $L_j \ge L_0 + j$, and
at $J^\star$ the potential has grown past $\log 2 + A\log 4$, which the
master law's lower half converts into a crossing.

\subsection{What the toll means}
Three readings, in decreasing order of how much of each is kernel-checked.

\emph{First: the price is a conservation law, not an optimization
target.} Theorem~\ref{thm:master} is two-sided. A cleverer tower of the
same shape cannot cross the budget more cheaply, because the decrement it
accumulates is its potential difference to within $5\%$, and the potential
difference is exactly what the crossing consumes. There is no slack to
find: the entire price of crossing $\log 2$ at design constant $A$ is
$4^{A}$ in doubly-logarithmic width, and both directions of that sentence
are theorems.

\emph{Second: the invoice has two axes, and a design chooses only the
axis.} At the flat threshold $\kappa = H/(A\log H)$, the multiplier $A$ is
a constant and the toll is a factor $4^{A}$ on the width
$c + \lambda = \log 2A + \log\log H$. At the original threshold
$\kappa = H/(\log H\log\log\log H)$, the role of $A$ is played by
$\log\log\log H$, which is not a constant; the same telescoping (with
potential $\log\log\log\log H$) then yields not a factor but a
\emph{power}: $(\log\log B)^{3} \le \log\log H_{J_{\min}} \le
(\log\log B)^{5}$
(\leaninline{Salt.Entropy.Chowla.tower\_loglog\_ge},
\leaninline{tower\_loglog\_le}). A power on the doubly-logarithmic axis is
a height on the scale axis, and it shows up as such: the budget floor of
the original road is the triple exponential
$\lceil e^{e^{e^{X}}}\rceil$ (\leaninline{budgetFloor}), while the flat
road's is the single exponential
$\lceil e^{\max(4X,\ 2\log A + 2)}\rceil$
(\leaninline{budgetFloorFlat}), with the design constant's own demand
$A \ge 2304\log 4/(\varepsilon^{6}\beta^{2})$
(\leaninline{budgetAFlat}) absorbing what the third logarithm used to
absorb. Two roads, one invoice, two axes. This paragraph's individual
claims are each kernel-checked. The sentence that fuses them --- that
\emph{every} threshold pays the same toll on one axis or the other --- is
proved below on the width axis, for every threshold schedule feeding this
step (Theorem~\ref{thm:shapetoll}); the further reading that the two axes
carry the same invoice remains a reading and not a theorem.

\emph{Third: the shape of the threshold does not change the price.} It is
natural to ask whether some other threshold shape evades the toll. Let
$\varphi : \mathbb{N} \to \mathbb{R}$ be an arbitrary \emph{schedule} --- a
different design constant at every level, chosen with full knowledge of the
tower --- and let the schedule tower and its telescoped decrement be
\[
H_0 = B, \qquad H_{j+1} = H_j\bigl\lfloor 2\varphi_j \log H_j\bigr\rfloor,
\qquad
S_J = \sum_{j<J} \frac{1}{2\varphi_j \log H_j}
\]
(\leaninline{Salt.Entropy.Chowla.chowlaTowerShape},
\leaninline{towerDropSumShape}). Fix a \emph{budget floor} $A_b \ge 1$ with
$\varphi_j \ge A_b$ at every level, and measure the width at that floor:
$c_b = \log 2A_b$, $\lambda_j = \log\log H_j$, $w_j = \log(c_b+\lambda_j)$
(\leaninline{towerLS, towerLamS, towerWS}).

\begin{theorem}[The shape-free toll]\label{thm:shapetoll}
Let $A_b \ge 1$, let $\varphi_j \ge A_b$ for every $j$, let
$B \ge 4\cdot10^{6}$, and let the flat floor hold at the budget constant,
$100(c_b + \lambda_0) \le L_0$. Then for \emph{every} $J$ at which the
schedule's decrement crosses the one-bit budget, $\log 2 < S_J$,
\[
\bigl(c_b + \lambda_0\bigr)\,4^{(20/21)A_b} \;\le\; c_b + \lambda_J .
\]
\leanname{Salt.Entropy.Chowla.towerShape\_width\_ge}
\end{theorem}

\begin{corollary}[The flat tower is extremal]\label{cor:extremal}
Take $A_b = A \ge 1$ and the hypotheses of Theorem~\ref{thm:shapetoll}.
Every schedule with budget floor $A$ multiplies the doubly-logarithmic
width by at least $4^{(20/21)A}$ when it crosses the budget; the constant
schedule $\varphi \equiv A$, which is the flat tower of
Definition~\ref{def:flattower}, multiplies it by at most
$\frac{21}{20}4^{A}$ at its minimal crossing length. The two bounds differ
by the factor $\frac{21}{20}\,4^{A/21}$.
\leanname{Salt.Entropy.Chowla.chowlaTowerShape\_const,
Salt.Entropy.Chowla.towerShape\_flat\_le}
\end{corollary}

Three observations make the schedule tower cost little beyond
\S\ref{sec:flat}. First, the per-step bracket \eqref{eq:stepbracket} is
proved over abstract reals (\leaninline{flatStep\_dw\_le},
\leaninline{flatStep\_dw\_ge}), so it is available at any step size.
Second, the floor price transfers in one direction only: $\lfloor\cdot\rfloor$
and $\log$ are monotone, so $\varphi_j \ge A_b$ gives
$\log\lfloor 2\varphi_j L_j\rfloor \ge c_b + \lambda_j - \frac{1}{500}$
(\leaninline{shapeStep\_logMul\_ge}), while the companion upper price
$\log\lfloor 2\varphi_j L_j\rfloor \le c_b + \lambda_j$ fails as soon as
$\varphi_j > A_b$. It is not needed: the necessity direction consumes only
the lower half. Removing it from the per-step law costs one monotonicity
step --- $w_{j+1} - w_j$ increases with the step, so truncating the step at
$\min(d_j,\, c_b + \lambda_j)$, which does satisfy both flat hypotheses,
only decreases it (\leaninline{shapeStep\_dw\_ge}). Third, the invariant
(H3) propagates for every schedule (\leaninline{towerShape\_inv}), and here
the schedule tower is the easier case: $\Delta\lambda_j \le d_j/L_j$ and
$100\,d_j/L_j \le d_j$ already at $L_j \ge 100$, so a larger multiplier
only helps.

Theorem~\ref{thm:shapetoll} quantifies over thresholds, not over step laws.
Its step is $H \mapsto H\lfloor 2\varphi\log H\rfloor$, the step that the
entropy decrement of \S\ref{sec:toll} funds; a decrement lemma whose
concatenation error is not of the form $\varepsilon^{2}/m$ would fund a
different step, and whether such a lemma exists is open. The exploratory
pass that preceded this theorem (nineteen probes, recorded in the
repository's ledger) also closed two escape routes that are not step-law
changes, cheap levels and $Q$-ary steps; those two closures are probe
records, not kernel-checked statements, and nothing above rests on them.

The practical corollary of the toll is one the campaign paid for three
times: an apparent wall in this method --- a triple-exponential floor, a
frame constant at a door, a window on the design constant --- is worth
re-reading as an invoice in the wrong coordinate before it is treated as a
height. In all three episodes the obstruction moved when the coordinates
changed. What did not move, in any of them, is $4^{A}$.
%
% TODO(three walls): \S\ref{sec:method} owes the three episodes by name,
% with dates and ledger entries, and owes the distinction between the two
% that dissolved and the one that did not.

% ============================================================================
\section{The flat tower}\label{sec:flat}
% ============================================================================

This section proves the per-step bracket \eqref{eq:stepbracket} on which
\S\ref{sec:toll} rests, and records the furniture that makes it available
at every level of the tower. Everything here is unconditional given (H1),
(H2), (H3).

\subsection{The tower does not stall}
The first thing to check about Definition~\ref{def:flattower} is that it
is a tower at all: if the multiplier were $1$ the sequence would be
constant and the decrement would never accumulate. Under (H1) and (H2),
$\lfloor 2A\log H\rfloor \ge 2$ for every $H \ge 4\cdot10^{6}$, since
$\log H \ge 15$ and $A \ge 1$
(\leaninline{Salt.Entropy.Chowla.flatMul\_ge\_two}). Hence every level
stays above the base floor,
$B \le H_j$ and $4\cdot10^{6}\le H_j$
(\leaninline{chowlaTowerFlat\_base\_ge}, \leaninline{chowlaTowerFlat\_base\_floor}),
and therefore $L_j \ge 15$ and $\lambda_j \ge 2$ at every level
(\leaninline{towerLF\_ge\_fifteen}, \leaninline{towerLamF\_ge\_two}). The
step in the glyph $L$ is exactly additive,
\[
L_{j+1} \;=\; L_j + d_j, \qquad d_j := \log\bigl\lfloor 2A L_j\bigr\rfloor ,
\]
\leaninline{Salt.Entropy.Chowla.towerLF\_succ}.

\subsection{The floor price}
The floor function in the multiplier is the only inexactness in the whole
construction, and it is priced on both sides.

\begin{lemma}[The floor price]\label{lem:floorprice}
Assume (H1), (H2). At every level, $d_j \le c + \lambda_j$. If moreover
the invariant $100(c+\lambda_j) \le L_j$ holds at that level, then
\[
c + \lambda_j - \tfrac{1}{500} \;\le\; d_j \;\le\; c + \lambda_j .
\]
\leanname{Salt.Entropy.Chowla.flatStep\_logMul\_le,
Salt.Entropy.Chowla.flatStep\_logMul\_ge}
\end{lemma}

The upper half is $\log\lfloor x\rfloor \le \log x$ with
$\log(2AL) = c + \lambda$ (\leaninline{flatLog\_prod}). The lower half is
$\log x - \log\lfloor x \rfloor \le (x - \lfloor x\rfloor)/\lfloor
x\rfloor \le 1/\lfloor x \rfloor$, and the invariant forces
$c + \lambda \ge 3$ (\leaninline{flatLevel\_s\_ge\_three}), hence
$L \ge 300$, hence $2AL \ge 600$ and $\lfloor 2AL\rfloor \ge 599$: the
floor costs at most $1/599 \le 1/500$ in the logarithm. That $1/500$ is
the only place where the discreteness of the construction is paid for, and
it is one of the three items inside the $21/20$ of
Theorem~\ref{thm:master}.

\subsection{The invariant}

\begin{proposition}[The flat floor propagates]\label{prop:inv}
Assume (H1), (H2) and the base hypothesis (H3). Then
$100(c + \lambda_j) \le L_j$ at every level $j$.
\leanname{Salt.Entropy.Chowla.towerFlat\_inv}
\end{proposition}

The induction is one line of arithmetic once Lemma~\ref{lem:floorprice} is
available: the step adds $d_j \ge 1$ to $L$ and at most $d_j/L_j \le
1/100$ to $\lambda$, so the slack $L - 100(c+\lambda)$ does not decrease.
The proposition is what allows every subsequent statement to assume the
invariant at an arbitrary level while hypothesizing it only at the base.

\subsection{The per-step bracket}
The heart of the toll is a statement about four real numbers and has
nothing to do with towers; the corpus states it that way, over abstract
reals, and then applies it at each level.

\begin{lemma}[The per-step bracket]\label{lem:step}
Let $L, L', d, c \in \mathbb{R}$ with $L' = L + d$, and set
$\lambda = \log L$, $\lambda' = \log L'$, $w = \log(c+\lambda)$,
$w' = \log(c+\lambda')$. Then:
\begin{enumerate}
\item[(i)] if $L > 0$, $c + \lambda > 0$, $0 \le d \le c + \lambda$, then
  $w' - w \le 1/L$;
\item[(ii)] if $c + \lambda \ge 3$, $100(c+\lambda) \le L$, and
  $c + \lambda - \tfrac{1}{500} \le d \le c + \lambda$, then
  $(20/21)/L \le w' - w$.
\end{enumerate}
\leanname{Salt.Entropy.Chowla.flatStep\_dw\_le,
Salt.Entropy.Chowla.flatStep\_dw\_ge}
\end{lemma}

Both halves are two applications of an elementary logarithm estimate, one
per glyph level. For (i) it is $\log x \le x - 1$ in the form
$\log b - \log a \le (b-a)/a$: first at the $L$ level,
$\lambda' - \lambda \le d/L$, then at the $c+\lambda$ level,
$w' - w \le (\lambda'-\lambda)/(c+\lambda)$; multiplying and using
$d \le c + \lambda$ gives $w' - w \le 1/L$ with no error term at all. For
(ii) it is the dual $1 - 1/x \le \log x$, giving
$d/L' \le \lambda' - \lambda$ and
$(\lambda'-\lambda)/(c+\lambda') \le w' - w$; the two second-order factors
$L' \le \frac{101}{100}L$ and $c + \lambda' \le (c+\lambda) + \frac{1}{100}$
(both from the invariant) and the floor price
$d \ge (c+\lambda) - \frac{1}{500}$ then give the constant. At
$c + \lambda \ge 3$ the constant this computation actually yields is
$0.986$; the lemma states $20/21 = 0.952\ldots$, which is what the
downstream arithmetic uses.

Applying Lemma~\ref{lem:step} at level $j$ with $d = d_j$, using
Proposition~\ref{prop:inv} for the invariant and
Lemma~\ref{lem:floorprice} for the floor price, gives
\eqref{eq:stepbracket} at every level
(\leaninline{Salt.Entropy.Chowla.towerStepF\_dw\_le},
\leaninline{towerStepF\_dw\_ge}), which is what \S\ref{sec:toll}'s proof
consumes.

\subsection{Reaching the crossing}
Proposition~\ref{prop:crossing} names a level at which the decrement has
crossed the budget. The estimate behind it is deliberately crude: since
$d_j \ge 1$ at every level, $L_j \ge L_0 + j$
(\leaninline{Salt.Entropy.Chowla.towerLF\_ge\_add}), so at
$J^\star = \lceil \exp(2(c+\lambda_0)4^{A})\rceil$ one has
$\lambda_{J^\star} \ge 2(c+\lambda_0)4^{A}$, hence
$w_{J^\star} - w_0 \ge \log 2 + A\log 4$, and the master law's lower half
turns that into $S_{J^\star} > \log 2$. The named level is superexponentially
larger than necessary --- the toll's own bracket \eqref{eq:bracket} shows
the crossing happens once the width has been multiplied by about $4^{A}$ ---
but it is explicit, which is what a regime builder needs: a Chowla regime
must exhibit a tower length, not merely know one exists.

The minimal length $J_{\min}$ is then the infimum of a nonempty set of
naturals (\leaninline{towerFlatJmin}), positive because $S_0 = 0 < \log 2$
(\leaninline{towerFlatJmin\_pos}), and every shorter tower stays under the
budget (\leaninline{towerDropSumFlat\_le\_log\_two\_of\_lt\_Jmin}) --- the
minimality fact that Theorem~\ref{thm:tollle} runs on.

\subsection{What the tower is for}
The flat tower was not built to prove the toll; it was built because the
original tower's third logarithm put a triple exponential in the budget
floor, and the toll is what the flat tower's audit produced. The regime
structure that consumes it carries the design constant $A$, the design law
$3.2A \le \log\log H_-$, a tower length, and the crossing
(\leaninline{Salt.Entropy.Chowla.ChowlaRegimeFlat}); the design law alone
discharges the flat floor (H3), the width export
(Corollary~\ref{cor:export}), and the height-1 budget floor. That
consumption is the subject of \S\ref{sec:pearl}.

% ============================================================================
\section{The pearl}\label{sec:pearl}
% ============================================================================

% ------------------------------------------------------------------
% STUB. Nothing in this section is drafted except the placeholder
% statement below and the hypothesis inventory, which are recorded here
% so that the final draft can be checked against them.
% ------------------------------------------------------------------

\noindent\textbf{TODO (\S\ref{sec:pearl} is a stub): this section awaits
BAND-WINDOW's final hypothesis list, so that the theorem is stated once,
in its final form.} What is pending, exactly: (a) whether the band-lane
hypothesis \leaninline{S16BandLaneCBoundedL} survives as a hypothesis or
is replaced by an $A$-windowed form; (b) which of the six numeric riders
($K_c$, $C_t$, $c_s$, $T_0$, $K_q$, $K_s$) have been discharged by then;
(c) the final form of the two carried predicates. Until (a)--(c) are
settled the statement below is a placeholder, and no prose in this section
should be written against it.

\begin{theorem}[The pearl; PLACEHOLDER STATEMENT]\label{thm:pearl}
For every $A_0 \ge 162$ there are constants
$\varepsilon, C_g, K_c, \delta_0, C_t, A, \beta, x_0, H_{\mathrm{opq}},
M_{\mathrm{fl}}, C_q, c_s, T_0, K_q, K_s, C$ with
$0 < \varepsilon$, $\tfrac{1}{500}\le\varepsilon$,
$A \ge 162$, $A \ge A_0$, together with a prefix of explicit sign and size
conditions on the constants and one arm-census conjunct pinning the
witness floor at the design base, such that under the eight named
hypotheses
\[
K_c \le 2^{539},\quad C_t \le 2^{23},\quad
x_0 \le e^{e^{3.2A}/10},\quad H_{\mathrm{opq}} \le \lceil e^{e^{3.2A}}\rceil,
\]
\[
e^{-100}\le c_s,\quad T_0 \le e^{e^{100}},\quad K_q \le e^{100},\quad
e^{-100}\le K_s,
\]
the following holds: for every $g : \mathbb{N}\to\mathbb{N}\to\mathbb{N}$
there is a Chowla regime $R$ with $R.\varepsilon = \varepsilon$, window
base $R.H_- $ pinned at the road's witness floor,
$g(R.H_+, R.\omega) \le R.x$, the width certificate
$\log\log R.H_+ \le 2e^{1.6A}$, and $3.2A \le \log\log R.H_-$, such that
the two carried predicates
\leaninline{S16CofactorSupply\_L\_gk} and
\leaninline{S16BaseScaleCap96\_L\_gk}
imply $\lnot\,\mathrm{logChowla2Fails}(R.\varepsilon, R.x, R.\omega)$.
\leanname{Salt.MR.logChowla2\_witnessed\_scale\_flat\_L\_v2
(Salt/MR/S16FlatFinal.lean)}
\end{theorem}

\noindent\textbf{TODO: prose for this section} --- the plain reading of
the conclusion (``at every design constant above the floor there is a
regime, arbitrarily deep, at which the correlation does not fail at the
witnessed scale''), the exact meaning of
\leaninline{logChowla2Fails} (namely
$\varepsilon\log\omega <
\bigl|\sum_{x/\omega < n \le x}\lambda(n)\lambda(n+1)/n\bigr|$,
\leaninline{Salt.Entropy.Chowla.logChowla2Fails}), the pinning of the base
at $\lceil e^{e^{3.2A}}\rceil$ (\leaninline{flatDesignBase}), the
arm-census conjunct of the prefix
($H_{\mathrm{opq}} \le \mathrm{flatDesignBase}(A) \Rightarrow
\mathrm{flatWitFloor}(\varepsilon,\beta,A,H_{\mathrm{opq}})
= \mathrm{flatDesignBase}(A)$), and how Corollary~\ref{cor:export}
supplies the regime's width conjunct. None of it is written yet.

% ============================================================================
\section{The classical limit}\label{sec:limit}
% ============================================================================

% ------------------------------------------------------------------
% HALF DRAFT. The containment argument and the quantifier-hoist finding
% are drafted; the final form of the corollary is TODO pending
% BAND-WINDOW (see \S\ref{sec:pearl}).
% ------------------------------------------------------------------

\subsection{The containment}
Two of Theorem~\ref{thm:pearl}'s eight hypotheses are windows rather than
bounds:
\[
x_0 \le e^{e^{3.2A}/10}
\qquad\text{and}\qquad
H_{\mathrm{opq}} \le \bigl\lceil e^{e^{3.2A}}\bigr\rceil .
\]
Each relates a constant produced by the theorem to a window that grows
with the design constant $A$, and $A$ is not bounded above: the theorem
produces, for every $A_0$, some $A \ge A_0$. That is the classical shape.
If the constants $x_0$ and $H_{\mathrm{opq}}$ do not themselves move with
$A_0$, then choosing
\[
A \;=\; \max\Bigl(\max(A_0, 162),\ \max\bigl(A_{\mathrm{budget}},\
\max(4x_0,\ H_{\mathrm{opq}})\bigr)\Bigr)
\]
discharges both windows outright, because the windows are double
exponentials and the constants are not: each falls to two applications of
$x + 1 \le e^{x}$. What is lost is effectivity --- the resulting $A$ is
bound by \texttt{Classical.choice} and no bound on it is claimed --- which
is precisely the caveat the analytic literature attaches to every
consequence of Siegel's theorem. The effective-conditional statement
contains the ineffective statement; the containment is the standard one,
and the corollary below is its instance.

\begin{corollary}[The Siegel-free form; ineffective in $A$]\label{cor:ineff}
Assume the band-lane hypothesis. For every $A_0$ there exist constants
$(\varepsilon, C_g, K_c, \delta_0, C_t, A, \beta, M_{\mathrm{fl}}, C_q,
c_s, T_0, K_q, K_s, C)$ with $A \ge 162$ and $A \ge A_0$ such that, under
the six riders
$K_c \le 2^{539}$, $C_t \le 2^{23}$, $e^{-100}\le c_s$,
$T_0 \le e^{e^{100}}$, $K_q\le e^{100}$, $e^{-100}\le K_s$, one has: for
every $g$ there is a Chowla regime $R$ with
$R.\varepsilon = \varepsilon$, base pinned at
$R.H_- = \lceil e^{e^{3.2A}}\rceil$, $g(R.H_+,R.\omega)\le R.x$,
$3.2A \le \log\log R.H_-$, $\log\log R.H_+ \le 2e^{1.6A}$, and such that
the two carried predicates imply
$\lnot\,\mathrm{logChowla2Fails}(R.\varepsilon,R.x,R.\omega)$.
No hypothesis of Siegel genre occurs: neither $x_0$ nor
$H_{\mathrm{opq}}$ appears in the statement.
\leanname{Salt.MR.logChowla2\_ineffective (Salt/MR/S16Uniform.lean)}
\end{corollary}

The base is exported as an \emph{equation}, $R.H_- = \lceil
e^{e^{3.2A}}\rceil$, not as a bound, and $A$ is unbounded: the regime is
pinned arbitrarily deep.

\subsection{The hoist}
The containment argument above has a hypothesis we underlined: \emph{if
the constants do not themselves move with $A_0$}. At the quantifier order
in which Theorem~\ref{thm:pearl} is stated, they may. The statement reads
\[
\forall A_0 \ge 162,\ \exists\,(\varepsilon, \ldots, x_0,
H_{\mathrm{opq}}, A, \ldots),\ \Phi(A_0, \varepsilon,\ldots),
\]
so every constant is chosen after $A_0$. Nothing in the statement asserts
that the witness for $A_0 = 162$ and the witness for $A_0 = 10^{6}$ share
their $x_0$; and a consumer cannot recover it, because two applications of
\texttt{Classical.choose} at two values of $A_0$ are choices from two
different propositions ($A_0$ occurs in $\Phi$, via $A_0 \le A$), so proof
irrelevance does not identify them. The classical limit is unavailable
until the order is corrected.

At the level of the proof term the constants are $A_0$-free: the caller's
$A_0$ enters the whole road at exactly one line, where the design constant
is chosen as $A := \max(A_0, A_{\mathrm{budget}})$, and
$\varepsilon$, $K_c$, $\delta_0$, $\beta$ and $H_{\mathrm{opq}}$ are all
minted strictly above it, while $x_0$ and the band constant come from a
hypothesis with no $A$ in it at all. Turning that reading of the proof
into a property of the \emph{statement} is the hoist: seven theorems along
the road were re-cut, each as the landed statement with $A$ moved out of
the existential and the corresponding $\forall A$ moved inward to the
point where the landed proof chose it, so that the road reads
\[
\exists\,(\text{constants}),\ \ldots \ \wedge\ \forall A,\ 162 \le A
\to A_{\mathrm{budget}} \le A \to \exists\,H_{\mathrm{cap}},\ (\text{the
same body}).
\]
Every body is the landed body verbatim.
\leanname{Salt.MR.flat\_head\_uniform, flat\_socket\_uniform,
flat\_doorL2\_uniform, flat\_road\_uniform, flat\_capstone\_uniform,\\
flat\_conditional\_uniform,
logChowla2\_witnessed\_scale\_flat\_L\_v2\_uniform}

We draw one methodological conclusion, and it is not about this theorem.
The step ``the constants are fixed, the window grows, so take the
parameter large'' is performed silently in analytic arguments hundreds of
times, and it is almost always right. Whether it is right in a given
instance is a question about quantifier order, which prose records only in
the reader's head, and which a formal statement records as data. Here the
step was blocked, the blockage was visible in the statement, and the
repair --- seven verbatim re-cuts --- was mechanical once identified.

\noindent\textbf{TODO (final form): this section is a half draft.} The
corollary above is stated as landed; BAND-WINDOW is expected to produce a
second corollary in which the band-lane hypothesis is itself $A$-windowed,
at which point (i) Corollary~\ref{cor:ineff} should be restated in that
final form, (ii) the sentence in \S\ref{sec:intro} listing the surviving
hypotheses must be updated to match, and (iii) the residual-hypothesis
table (Table~\ref{tab:conditional}) must be regenerated. The paragraph on
the hoist is stable and does not depend on that outcome.

% ============================================================================
\section{Method}\label{sec:method}
% ============================================================================

% ------------------------------------------------------------------
% STUB. Contents listed, prose not written. The full treatment is the
% companion paper \cite{companion}; this section is the part of it that
% this paper's results require.
% ------------------------------------------------------------------

\noindent\textbf{TODO (\S\ref{sec:method} is a stub).} Intended content,
in the order it should be written:

\begin{enumerate}
\item \emph{The design-block and refuter protocol.} Statements are frozen
by a design block before any proof is attempted; a freeze is then attacked
by independent refuters with a standing default of \emph{refuted}, each
carrying assigned kill-checks; executors are forbidden to alter a frozen
statement and are required to stop and report on discovering one
unprovable. To be written with the three refuter verdicts on the pearl
(the non-vacuity pass, the satisfiability pass, and the delta pass over
the final supply chain) as the worked example, including the two defects
each found and how they were repaired at the statement level.

\item \emph{The discovery mode.} How the toll was found: not transcribed
from the literature, but produced by auditing what a formalized tower
actually charges, at a point in the campaign where three different
apparent walls were open at once. The claim to be made carefully here is
about the mode --- a two-sided law came out of a demand audit that a prose
development has no reason to perform --- and not about novelty of the
underlying computation.

\item \emph{The three walls, by name.} The triple-exponential budget
floor; the frame constant at the door; the window on the design constant.
Two dissolved on re-reading and one did not; the section owes which is
which, with dates and ledger entries, and owes the general lesson stated
in \S\ref{sec:toll}: an apparent wall in this method is worth re-reading
as an invoice in the wrong coordinate before it is treated as a height.

\item \emph{Walls that stayed walls.} The section must state plainly what
this method did not get past, including the items still carried as
hypotheses in \S\ref{sec:pearl} and the two constants whose traces are
open. A method section that lists only its successes is not evidence.

\item \emph{The error ledger.} The public, append-only record of every
design error caught, numbered; the count at the time of writing, and the
measurement that goes with it (no design error first discovered by the
kernel).
\end{enumerate}

% ============================================================================
\section*{Acknowledgments}
% ============================================================================
This work was conducted entirely in the author's personal capacity, on
personal time and personal equipment, and used no confidential or
proprietary information of any employer. The views, findings, and
conclusions expressed here are the author's own and do not represent
the views, positions, or endorsement of the author's employer or of any
other organization. The collaboration with Claude (Anthropic's Fable~5
and Opus models) is described in \S\ref{sec:method} and recorded
commit-by-commit in the repository; it likewise implies no position or
endorsement on the part of Anthropic. The author directed and ratified
every design decision and bears sole responsibility for the claims. The
formalization builds on \textsf{mathlib}, and this work stands on the
shoulders of that community.

% ============================================================================
\appendix
\section{The formal-statement index}\label{app:index}
% ============================================================================
% At release: the axiom column is regenerated from the All.lean
% #audit_axioms blocks, exactly as in the flagship. ``The three'' below
% records the corpus-wide audit at the sealed state (build 9674/9675,
% exit 0); rows marked (stub) await \S\ref{sec:pearl}'s final statement.

\begin{table}[h]
\centering\small
\begin{tabular}{lll}
\toprule
Statement & Declaration & Axioms \\
\midrule
\multicolumn{3}{l}{\emph{The toll} (\S\ref{sec:toll}, \S\ref{sec:flat});
file \texttt{Salt/Entropy/Chowla/TowerFlat.lean} unless noted} \\
\midrule
Def~\ref{def:flattower} & \leaninline{chowlaTowerFlat} & --- \\
Def~\ref{def:flattower} & \leaninline{towerDropSumFlat} & --- \\
Def~\ref{def:potential} & \leaninline{flatC, towerLF, towerLamF, towerWF} & --- \\
Thm~\ref{thm:master} (lower) & \leaninline{towerDropSumFlat\_ge\_log\_ratio} & the three \\
Thm~\ref{thm:master} (upper) & \leaninline{towerDropSumFlat\_le\_log\_ratio\_mul} & the three \\
Thm~\ref{thm:tollge} & \leaninline{towerFlat\_width\_ge} & the three \\
Thm~\ref{thm:tollle} & \leaninline{towerFlat\_width\_le} & the three \\
Prop~\ref{prop:crossing} & \leaninline{towerFlatJstar} (def) & --- \\
Prop~\ref{prop:crossing} & \leaninline{towerFlat\_crossing\_at\_Jstar} & the three \\
Prop~\ref{prop:crossing} & \leaninline{towerFlat\_crossing\_exists} & the three \\
Prop~\ref{prop:crossing} & \leaninline{towerFlatJmin} (def), \leaninline{towerFlatJmin\_spec} & the three \\
Prop~\ref{prop:crossing} & \leaninline{towerFlatJmin\_pos} & the three \\
\S\ref{sec:toll} minimality & \leaninline{towerDropSumFlat\_le\_log\_two\_of\_lt\_Jmin} & the three \\
Cor~\ref{cor:export} & \leaninline{towerFlat\_width\_export}
  (\texttt{TowerFlatExport.lean}) & the three \\
Cor~\ref{cor:export} engine & \leaninline{flatWidth\_coef, flatWidth\_engine\_at,
  flatWidth\_engine} & the three \\
(H3) is a theorem & \leaninline{flatFloor\_of\_design}
  (\texttt{TowerFlatRegime.lean}) & the three \\
Lem~\ref{lem:floorprice} & \leaninline{flatStep\_logMul\_le, flatStep\_logMul\_ge} & the three \\
Prop~\ref{prop:inv} & \leaninline{towerFlat\_inv} & the three \\
Lem~\ref{lem:step} & \leaninline{flatStep\_dw\_le, flatStep\_dw\_ge} & the three \\
\eqref{eq:stepbracket} at level $j$ & \leaninline{towerStepF\_dw\_le, towerStepF\_dw\_ge} & the three \\
\S\ref{sec:flat} tower floor & \leaninline{flatMul\_ge\_two,
  chowlaTowerFlat\_base\_floor} & the three \\
\S\ref{sec:flat} glyph step & \leaninline{towerLF\_succ, towerLF\_ge\_add} & the three \\
\S\ref{sec:flat} level bounds & \leaninline{towerLF\_ge\_fifteen,
  towerLamF\_ge\_two,} & the three \\
 & \leaninline{flatLevel\_s\_ge\_three} & \\
\S\ref{sec:toll} threshold funding & \leaninline{flat\_hkey} & the three \\
\S\ref{sec:toll} budget endpoints & \leaninline{entropy\_per\_symbol\_le,}
  & the three \\
 & \leaninline{entropy\_nonneg\_per\_symbol} (\texttt{Endpoints.lean}) & \\
\S\ref{sec:toll} telescoping & \leaninline{tower\_telescopeFlat}
  (\texttt{DecrementFlat.lean}) & the three \\
\S\ref{sec:toll} the other axis & \leaninline{tower\_loglog\_ge,
  tower\_loglog\_le} (\texttt{TowerExport}) & the three \\
\S\ref{sec:toll} the two floors & \leaninline{budgetFloor}
  (\texttt{BudgetCore}), \leaninline{budgetFloorFlat,} & --- \\
 & \leaninline{budgetAFlat} (\texttt{BudgetFlat.lean}) & \\
\midrule
\multicolumn{3}{l}{\emph{The shape-free toll} (\S\ref{sec:toll});
file \texttt{Salt/Entropy/Chowla/TowerShape.lean}} \\
\midrule
Thm~\ref{thm:shapetoll} schedule & \leaninline{chowlaTowerShape,
  towerDropSumShape} (defs) & --- \\
Thm~\ref{thm:shapetoll} glyphs & \leaninline{towerLS, towerLamS, towerWS,
  towerMulS} (defs) & --- \\
Thm~\ref{thm:shapetoll} & \leaninline{towerShape\_width\_ge} & the three \\
Thm~\ref{thm:shapetoll} telescope & \leaninline{towerShape\_dropSum\_le} & the three \\
Thm~\ref{thm:shapetoll} invariant & \leaninline{towerShape\_inv} & the three \\
Thm~\ref{thm:shapetoll} floor price & \leaninline{shapeStep\_logMul\_ge\_abs,
  shapeStep\_logMul\_ge} & the three \\
Thm~\ref{thm:shapetoll} per-step & \leaninline{shapeStep\_dw\_ge,
  towerStepS\_dw\_ge} & the three \\
Cor~\ref{cor:extremal} tie & \leaninline{chowlaTowerShape\_const, towerLS\_const,}
  & the three \\
 & \leaninline{towerLamS\_const, towerWS\_const, towerDropSumShape\_const} & \\
Cor~\ref{cor:extremal} & \leaninline{towerShape\_flat\_le} & the three \\
\S\ref{sec:toll} schedule floor & \leaninline{shapeMul\_ge\_two,
  chowlaTowerShape\_base\_floor,} & the three \\
 & \leaninline{towerLS\_ge\_fifteen, towerLamS\_ge\_two, towerLS\_succ} & \\
\midrule
\multicolumn{3}{l}{\emph{The pearl} (\S\ref{sec:pearl}, \S\ref{sec:limit})
--- statements stubbed, declarations landed} \\
\midrule
Thm~\ref{thm:pearl} (stub) &
  \leaninline{Salt.MR.logChowla2\_witnessed\_scale\_flat\_L\_v2} & the three \\
Cor~\ref{cor:ineff} &
  \leaninline{Salt.MR.logChowla2\_ineffective} & the three \\
\S\ref{sec:limit} hoist (7) &
  \leaninline{Salt.MR.flat\_head\_uniform \ldots} & the three \\
 & \leaninline{logChowla2\_witnessed\_scale\_flat\_L\_v2\_uniform} & \\
\S\ref{sec:pearl} target &
  \leaninline{Salt.Entropy.Chowla.logChowla2Fails} (def) & --- \\
\bottomrule
\end{tabular}
\caption{``The three'' $=$ \leaninline{propext, Classical.choice,
Quot.sound}. Toll rows are complete; pearl rows carry landed declarations
against statements that \S\ref{sec:pearl} has not yet fixed. Full index
generated at release.}
\label{tab:index}
\end{table}

\begin{table}[h]
\centering\small
\begin{tabular}{lll}
\toprule
Conditional statement & Residual hypothesis & Status \\
\midrule
Thm~\ref{thm:pearl} & band-lane constant; 8 numeric riders; & \S\ref{sec:pearl}
  (stub) \\
 & 2 carried predicates & \\
Cor~\ref{cor:ineff} & band-lane constant; 6 numeric riders; & landed \\
 & 2 carried predicates & \\
\bottomrule
\end{tabular}
\caption{Every conditional in this paper, with its exact named residual.
Nothing in \S\ref{sec:toll} or \S\ref{sec:flat} is conditional: the toll
is unconditional given its three stated hypotheses (H1)--(H3), which are
hypotheses on the tower's parameters, not on the theory.
\textbf{TODO: regenerate when \S\ref{sec:pearl} is drafted.}}
\label{tab:conditional}
\end{table}

\begin{thebibliography}{99}
\bibitem{TaoChowla} T.~Tao, \emph{The logarithmically averaged Chowla and
Elliott conjectures for two-point correlations}, Forum Math.\ Pi
\textbf{4} (2016).
\bibitem{MR2016} K.~Matom\"aki and M.~Radziwi\l\l, \emph{Multiplicative
functions in short intervals}, Ann.\ of Math.\ \textbf{183} (2016).
\bibitem{MRT} K.~Matom\"aki, M.~Radziwi\l\l, and T.~Tao, \emph{An
averaged form of Chowla's conjecture}, Algebra Number Theory \textbf{9}
(2015).
\bibitem{Chowla} S.~Chowla, \emph{The Riemann hypothesis and Hilbert's
tenth problem}, Gordon and Breach, 1965.
\bibitem{mathlib} The mathlib community, \emph{The Lean mathematical
library}, CPP 2020.
\bibitem{fulcrum} J.~Hickey, \emph{Twin primes and Siegel zeros: a
fulcrum}, in preparation.
\bibitem{companion} J.~Hickey, \emph{The Salt method}, in preparation.
\end{thebibliography}

\end{document}
